\subsection{Stalks and the absolute sign}
Two cylindrical completions of the same noncompact smooth graph need not give the same object in a stopped category. We compute the class of $B_P$ by testing it at the strata of the angle torus. Write $\mu_{\rm FZ}$ for the covariant cotangent microlocalization with primitive $-y\,dm$ in \cite[Theorem~5.1]{FZ}.

\begin{lemma}\label{lem:stalkprobe}
With the costandard grading on its inner logarithmic graph, the class of $\mu_{\rm FZ}^{-1}B_P$ is $[\kappa\mathcal O(P)]$. In the orientation of this grading its rightward class is $[C_P]$, and its endpoint allowed class is the corresponding signed positive displacement.
\end{lemma}

\begin{proof}
Use the same subsets in coordinates $(m,y)$ with the GPS primitive $+y\,dm$ for the calculation. Changing the sign of the primitive reverses the category, by reversing boundary motion and the order of inputs in the disks defining its operations; use the vertically polarized decorations under this identification. The opposite GPS equivalence thus gives the stated covariant convention. We first compare the classes up to sign.

Make a finite affine triangulation $\mathcal S$ of the angle torus refining all projected faces of $P$ and all annihilator subtori. For instance take the periodic hyperplanes $\langle m,v\rangle\in\mathbb Z$ of the rays, add a fine periodic coordinate grid to embed local stars, and barycentrically subdivide on the cover. The whole conormal stop of $\mathcal S$ contains the fan stop. The small displaced normal link also misses this larger stop. Otherwise a segment of length $e$ in direction $\nu$ from a face of $P$ to a simplex on the cover would be perpendicular to both affine hulls. Its length would be the distance between them. Only finitely many lifted simplices meet a fixed neighbourhood of $P$; choose $e$ smaller than every strictly positive such distance. The uniform estimates above allow $B$ large enough to avoid this stop on the actual cylindrical disk.

For the $+y\,dm$ primitive a cofinal positive wrapping of this disk is obtained by decreasing $e$ to zero while increasing $B=B(e)$ to infinity. Indeed its angular support function at infinity is $k_e(\nu)=g(B(e)\nu)/B(e)-e$, with $\partial_e k_e<0$ since $\partial_B(g(B\nu)/B)=b_P(dg(B\nu))/B^2>0$ and $B(e)$ is decreasing. Thus decreasing $e$ gives positive contact motion. Take $B(e)$ large enough on successive compact intervals that the boundary misses the triangulation stop throughout. The links eventually leave every compact subset of its complement by their uniform convergence to the undisplaced normal relation. Cofinality follows from \cite[Lemma~5.1]{GPSmicro}.

For a point $q$ in the relative interior of a simplex take a rounded outward conormal to a small ball at $q$. In a flat chart use the graph $m=q+\delta d\rho(y)$, with $\rho$ radial convex and equal to $|y|$ plus a constant at infinity, and take the grading of the cotangent fibre. Its boundary avoids the triangulation stop for sufficiently small $\delta$ by the same distance argument. Pairing with this test brane in the second slot computes stalk Euler characteristic. Indeed on the open-star generators the Euler number is one when the star contains $q$ and zero otherwise by \cite[\S5.7, especially Lemma~5.21, and Theorem~5.34]{GPSmicro}. These generators give the finite incidence category. Filtration by open simplices, or the triangular matrix of star indicators, shows that stalk Euler numbers on the strata determine the integral class there.

The moving links of the cofinal wrapping eventually miss the fixed test end throughout, so continuation stabilizes and the Euler number is the signed intersection with a sufficiently far wrapped disk. The prisms of the preceding subsection permit its replacement for this intersection by the undisplaced outward normal cycle. In fact the test end is disjoint from the raw normal link, a compact subset of the triangulation stop. For small $e$ and then large conification radii all link traces miss that test end throughout the interpolation, including removal of the displacement. Interior perturbations do not change the intersection number.

Sum over the finitely many relevant lifts of $q$. An exterior lift gives zero for $\delta$ below its distance to $P$. At an interior lift shrinking the ball gives the single transverse intersection of the fibre with the open zero section, of constant sign. At a boundary lift choose $z\in P^\circ$, set $v=q-z$, and move the test centre to $q+uv$ for small $u\ge0$. A boundary coincidence would have $m=q+uv+\delta\nu$ for an outward unit normal $\nu$ to $P$ at $m$. All faces meeting a sufficiently small neighbourhood of $q$ contain $q$. Thus $\langle\nu,m-q\rangle=0$, whereas $\langle\nu,v\rangle>0$. This contradicts $0=u\langle\nu,v\rangle+\delta$. When the centre reaches an exterior point one can choose $\delta$ smaller than its distance to $P$. This proves the vanishing of each boundary contribution by an intersection homotopy.

Up to sign the stalk function is therefore the number of interior lifts into $P^\circ$. The sheaf class is $\pm[\pi_!j_!\omega_{P^\circ}]$ for the inclusion $j$ of the open polytope and angular projection $\pi$; its dualizing shift has a single fixed parity. This sheaf already has microsupport in the negative fan cone. When one forgets the added triangulation stops, the sheaf functor is the left adjoint of inclusion of the fan microsupport category and fixes this sheaf. The brane functor takes our disk to itself. In the fan category its class is consequently that of the prescribed costandard brane $A_P$ of \cite[\S5.2, Proposition~5.2 and Corollary~5.3]{FZ}, up to sign.

Give $B_P$ the grading of $A_P$ on the unmodified inner graph. Both meet the compact zero section in the same transverse point at $y=0$ with the same graded orientation. Its Euler pairing is a unit for both, fixing the categorical sign. For the geometric sign instead use the compact angle torus at small constant momentum $y_0$ in the open normal cone of a vertex. It pairs with rightward relative classes by choosing the threshold above its compact image. It meets each of the two branes once, at the same point with the same orientation, and meets the raw normal cycle over the chosen vertex with multiplicity one. The physical class of $A_P$ is $[C_P]$ by its own cycle comparison \cite[\S5.2]{FZ}. Hence the geometric sign for $B_P$ gives $[C_P]$ as well. Its prism in the allowed group used this very same orientation. This proves all three assertions using the explicit cylindrical disk.
\end{proof}

\subsection{The marked collection}\label{sec:markedtheorem}
Let $I_a$ denote the standard rightward-cycle-to-$K$ isomorphism at the uncut regular ray, normalized by the positive real cycle for $[\mathcal O_X]$. Equation~\eqref{eq:cycledual} states that $I_a[C_P]=[\mathcal O(P)^\vee]$. On the opposite of the category $\mathcal C_a$ of the compact fibration at $a$, define the complex-linear equivalence
\[
       V=D_X\circ\kappa^{-1}\circ\mu_{\rm FZ}^{-1}\circ T_*
       :\mathcal C_a^{\mathrm{op}}\longrightarrow\mathrm{Perf}(X),
\]
where the contravariant composite is written on objects and regarded as covariant out of the opposite. We claim, identifying opposite Grothendieck groups by the class of an object, that
\begin{equation}\label{eq:raymarking}
                         K_0(V)=I_a\,J_a\,c_a .
\end{equation}
For $x_P=[T_*^{-1}B_P]$ both evaluations are explicit:
\[
 \begin{aligned}
      K_0(V)(x_P)&=[\mathcal O(P)^\vee],\\
      J_a c_a(x_P)&=[C_P],\qquad I_a[C_P]=[\mathcal O(P)^\vee].
 \end{aligned}
\]
The first equality is the stalk calculation in Lemma~\ref{lem:stalkprobe}. The second uses the brane fundamental class of Lemma~\ref{lem:braneclass} and the positive displacement of Lemma~\ref{lem:positivepush}, with the same orientation. Take $P$ in a finite list whose ample line bundles span integral $K$-theory, as in Proposition~\ref{prop:cycleduality}. Then $x_P$ span by the endpoint equivalences, proving \eqref{eq:raymarking}.

All comparisons in the two evaluations can be made on the same compact domains. For the finite list first choose displacements and $D_-,D$, then $\beta$. Take $L'\ge A_-$ above the regular-page thresholds for both paths and $L$ beyond all disk traces before fixing the common $\kappa$. At $*$ reduce the page level to $A_-$ through its regular collar. Singular interior fibres on the real coefficient path leave the boundary regular.

At $a$ the preferred disk is a thimble of the genuine uncut Laurent polynomial. Adding base area does not change its fibre connection. An interpolation of positive Legendre metrics on the compact part transports its collapsing cap and finite regular segments in a compact cobordism, so it defines the same relative class. The bend of the preferred ray towards the positive comparison page can be taken above the threshold in \eqref{eq:relativejunction}, using our small initial phase and its lift. Thus $J_a c_a$ of this disk is its usual rightward thimble. The small toric mirror period comparison takes this cycle to the asymptotic column of its own critical point, up to orientation \cite[\S6.2]{FZ}, \cite[proof of Theorem~6.4]{GI}, \cite[Theorem~7.5]{GHIKLS}. More precisely the integrals with the Brieskorn insertions give a vector solution on a sector of opening greater than $\pi$, with the leading evaluation at the given Morse point \cite[Proposition~6.18]{IriGlobal}. In the convention $e^{F/z}$ there, replace $z$ by $-z$ and the phase by $\varphi+\pi$ with the corresponding half-turn lift to obtain our $e^{-W/z}$ convention. Sectorial uniqueness identifies the full column. A leading term on a single ray would permit exponentially small additions. This use of the period comparison is entirely at $a$.

\begin{theorem}\label{thm:markedcollection}
Let $X$ be a smooth projective complex toric Fano variety and let $G$ be a finite group preserving its fan. There is a $G$-fixed semisimple point $\tau$ on the convergent large-radius branch such that, for every admissible lifted phase $\varphi$, the primitive idempotents $s$ label a full exceptional collection $(E_s)$ in $D^b\mathrm{Coh}(X)$ with the following properties. In decreasing order of $\operatorname{Im}(e^{-i\varphi}u_s)$, its individual Gamma sections are the corresponding normalized asymptotic columns, with orbit-compatible signs. Inverse toric pullback carries $E_s$ to $E_{gs}$ up to isomorphism. For $s\ne s'$ with $u_s=u_{s'}$ one has $\mathrm{RHom}(E_s,E_{s'})=0$. In particular the subcategories at the distinct Euler values give a $G$-stable semiorthogonal decomposition; the subcategory at a value of multiplicity $m$ is equivalent to $\mathrm{Perf}(\C)^m$.
\end{theorem}

\begin{proof}
For the initial lifted phase, \eqref{eq:raymarking} verifies the premise of Proposition~\ref{prop:markreduction}. Continue $V$ by that proposition to obtain the unsigned classes of the individual $D_s$ at $\tau$. On this same opposite category $D_X U$ gives the permuted reference collection in decreasing-height order by Theorem~\ref{thm:coherentcollection}. Apply Proposition~\ref{prop:ordinarymarking}, using the integral orbit-normalized basis of Proposition~\ref{thm:main}. It yields the stated permutation of coherent objects and signs constant on orbits. Shifts on the negative orbits give the exact prescribed classes. Exceptionality and the two Hom vanishings at equal values follow from the original disk calculation. Changing the lifted phase uses the block mutations of Section~\ref{sec:blocktransport}; its continuation preserves the individual labels and acts on the Gamma basis by the same Euler mutations. For a point use the structure sheaf. The parameter and initial choices can be made for the full fan group, then restricted to any subgroup.
\end{proof}

The signed collection need not coincide term by term with that constructed using the covariant functor \eqref{eq:coherentendpoint}. The comparison uses both functors on specified geometric disks and supplies an ordinary collection with the two required properties. On a stratum with constant equality partition the phase transport of Section~\ref{sec:blocktransport} applies to its blocks. Orbit decomposition on one fibre and connected components of the Euler cover over a varying base are different operations; we determine the latter next.
